%! Conic Sections — Unit 8, Lesson 8.5 — Lesson Plan
\documentclass[10pt]{article}
\usepackage{precalculus-article}
\usepackage{precalculus-boxes}
\usepackage{graphicx}
\graphicspath{{images/}}
\newcommand{\TallMath}[1]{$\displaystyle #1\rule[-1.4em]{0pt}{3.2em}$}
\newcommand{\CourseName}{Precalculus}
\newcommand{\MeetingLength}{55 minutes}
\newcommand{\UnitNumberName}{Unit 8: Parametric Functions, Vectors, and Matrices \quad}
\newcommand{\LessonNumberName}{Lesson 8.5: Conic Sections}

\begin{document}

%----------------------------------------------------------------------
% Title Block
%----------------------------------------------------------------------
\begin{center}
  {\Large\bfseries \CourseName} \\[4pt]
  {\small\bfseries \UnitNumberName \LessonNumberName \quad (3 instructional periods, \MeetingLength\ each)}
\end{center}

%----------------------------------------------------------------------
% Primary Objective
%----------------------------------------------------------------------
\begin{tcolorbox}[colback=sky, colframe=navy, boxrule=0.9pt, arc=2mm,
    left=2mm, right=2mm, top=1.5mm, bottom=1.5mm]
  \textbf{Primary Objective:} Students will write, interpret, and convert
  analytic representations of parabolas, ellipses, circles, and hyperbolas
  with horizontal or vertical lines of symmetry, identifying center or vertex,
  axes radii, and asymptotes from standard form equations.
  \textit{(Big Idea: Functions)}
\end{tcolorbox}

\vspace{0.08in}

%----------------------------------------------------------------------
% Priority Ideas & Skills
%----------------------------------------------------------------------
\begin{skillbox}[Priority Ideas \& Skills]{goldbox}
\begin{tabularx}{\linewidth}{@{} p{0.46\linewidth} X @{}}
\textbf{AP Skills} & \textbf{Key Understandings (EKs)} \\
\midrule
\textbf{1.B} \textit{Express in equivalent forms} ---
  Rewrite a general second-degree equation (e.g.\ $4x^2 - 9y^2 = 36$) in
  standard conic form by dividing through or completing the square, then
  identify the conic type and its key measurements.
&
\textbf{EK 4.6.A.1}: A parabola with vertex $(h,k)$ can, if $a \ne 0$, be
represented analytically as $x - h = a(y - k)^2$ if it opens left or right,
or as $y - k = a(x - h)^2$ if it opens up or down. \\[4pt]

\textbf{2.A} \textit{Identify information from representations} ---
  Extract the center, vertex, semi-axes, and asymptotes directly from a
  standard-form equation without graphing.
&
\textbf{EK 4.6.A.2}: An ellipse centered at $(h,k)$ with horizontal radius $a$
and vertical radius $b$ is represented as
$\dfrac{(x-h)^2}{a^2} + \dfrac{(y-k)^2}{b^2} = 1$.
A circle is a special case where $a = b$. \\[4pt]

\textbf{2.B} \textit{Construct equivalent representations} ---
  Given a verbal or geometric description of a conic (center or vertex plus key
  measurements), write the standard-form equation.
&
\textbf{EK 4.6.A.3}: A hyperbola centered at $(h,k)$ is
$\dfrac{(x-h)^2}{a^2} - \dfrac{(y-k)^2}{b^2} = 1$ (opens left/right)
or $\dfrac{(y-k)^2}{b^2} - \dfrac{(x-h)^2}{a^2} = 1$ (opens up/down),
with asymptotes $y - k = \pm\dfrac{b}{a}(x - h)$. \\
\end{tabularx}
\end{skillbox}

\vspace{0.08in}

%----------------------------------------------------------------------
% Vocabulary
%----------------------------------------------------------------------
\begin{skillbox}[{Vocabulary, Concepts, \& Theorems}]{greenbox}
\begin{tabularx}{\linewidth}{@{} p{0.26\linewidth} X @{}}
\textbf{Term} & \textbf{Definition / Note} \\
\midrule
conic section &
  A curve formed by the intersection of a plane and a double cone:
  parabola, ellipse (including circles), or hyperbola. \\[4pt]
parabola &
  The set of points equidistant from a fixed focus and a fixed directrix;
  vertex $(h,k)$ in standard form. \\[4pt]
vertex &
  The turning point of a parabola; the point $(h,k)$
  in $y - k = a(x-h)^2$ or $x - h = a(y-k)^2$. \\[4pt]
ellipse &
  The set of points for which the sum of distances to two foci is constant;
  has a center and two perpendicular axes. \\[4pt]
center &
  The point $(h,k)$ at the center of an ellipse or hyperbola. \\[4pt]
semi-major / semi-minor axis &
  The larger ($a$) and smaller ($b$) half-axes of an ellipse. \\[4pt]
circle &
  Special ellipse with $a = b = r$; equation
  \TallMath{(x-h)^2 + (y-k)^2 = r^2}. \\[4pt]
hyperbola &
  The set of points for which the \emph{difference} of distances to two foci
  is constant; has two branches and two asymptote lines. \\[4pt]
asymptote &
  Lines the hyperbola approaches but never reaches:
  \TallMath{y - k = \pm\tfrac{b}{a}(x - h)}. \\
\end{tabularx}
\end{skillbox}

\vspace{0.08in}

%----------------------------------------------------------------------
% Activate Prior Knowledge
%----------------------------------------------------------------------
\begin{skillbox}[Activate Prior Knowledge \& Spiral Review]{sky}
\begin{tabularx}{\linewidth}{@{}X c@{}}
\begin{minipage}[t]{0.96\linewidth}\vspace{0pt}
\textbf{Spiral topics activated during Warm-Up (first 8 minutes):}
\begin{itemize}[leftmargin=1.4em, itemsep=2pt]
  \item \textbf{Completing the square}: Rewrite $x^2 - 6x$ as $(x-3)^2 - 9$
    --- the essential algebraic tool for converting general equations to
    standard conic form.
  \item \textbf{Equation of a circle}: Recall $(x-h)^2 + (y-k)^2 = r^2$ as a
    familiar instance of the ellipse standard form, grounding new content in
    prior knowledge.
  \item \textbf{Reading standard form}: Identify center and radii from a
    given ellipse equation --- directly previews extracting features from
    ellipse and hyperbola equations.
\end{itemize}
\end{minipage}
&
\end{tabularx}
\end{skillbox}

\vspace{0.08in}

%----------------------------------------------------------------------
% Hook
%----------------------------------------------------------------------
\begin{skillbox}[Hook --- Entry Event]{sky}
\textbf{Scenario:} ``Planetary orbits are ellipses. Telescope mirrors are
parabolas. Navigation systems like LORAN use hyperbolas. What do all these
shapes have in common, and how do we write equations that pin them down
precisely?''

\textbf{Discussion prompts:}
\begin{enumerate}[leftmargin=1.4em, itemsep=3pt]
  \item \textit{``A satellite dish focuses signals at one point. What shape is
    the dish, and why does that shape work?''}
  \item \textit{``Earth's orbit is an ellipse with the Sun at one focus. How
    is an ellipse different from a circle?''}
  \item \textit{``Two radio stations broadcast simultaneously. A ship measures
    the difference in arrival times. The set of positions with a fixed
    difference in distances to the two stations is what shape?''}
\end{enumerate}

\textbf{Key tension to surface:} All three conics share elegant standard-form
equations. Mastering those equations lets us describe physical curves precisely.

\textbf{Transition:} ``Over three periods we master the standard form of each
conic, identify key features, and write equations from descriptions.''
\end{skillbox}

\vspace{0.08in}

%----------------------------------------------------------------------
% Lesson — Period 1: Parabolas
%----------------------------------------------------------------------
\begin{skillbox}[Lesson --- Period 1: Parabolas]{sky}
\begin{multicols}{2}

\textbf{Step 1 --- Two Orientations} (10 min)

Present both standard forms:
\begin{itemize}[leftmargin=1.2em, itemsep=2pt]
  \item \textbf{Vertical axis:} $y - k = a(x - h)^2$
    ($a > 0$ opens up; $a < 0$ opens down)
  \item \textbf{Horizontal axis:} $x - h = a(y - k)^2$
    ($a > 0$ opens right; $a < 0$ opens left)
\end{itemize}
Emphasize: the \emph{squared variable} tells you the axis of symmetry.

\columnbreak

\textbf{Step 2 --- Reading the Equation} (20 min)

Model: $y + 3 = 2(x - 1)^2$.
Rewrite as $y - (-3) = 2(x-1)^2$.
Vertex $(1, -3)$; $a = 2 > 0$ $\Rightarrow$ opens \emph{up}.

Student practice:
\begin{itemize}[leftmargin=1.2em, itemsep=2pt]
  \item $y - 5 = -(x + 2)^2$ --- vertex and direction?
  \item $x - 4 = -\tfrac{1}{2}(y - 1)^2$ --- vertex and direction?
\end{itemize}

Ask: \textit{``How does the sign of $a$ tell you direction without
graphing?''}

\end{multicols}

\smallskip
\textbf{Pacing:} Reserve final 15 min for students to start Tier R portion of activity.
\end{skillbox}

\vspace{0.08in}

%----------------------------------------------------------------------
% Lesson — Period 2: Ellipses and Circles
%----------------------------------------------------------------------
\begin{skillbox}[Lesson (cont.) --- Period 2: Ellipses and Circles]{sky}
\begin{multicols}{2}

\textbf{Step 3 --- Ellipse Standard Form} (15 min)

Standard form: $\dfrac{(x-h)^2}{a^2} + \dfrac{(y-k)^2}{b^2} = 1$

Example: $\dfrac{(x+1)^2}{16} + \dfrac{(y-3)^2}{9} = 1$
\begin{itemize}[leftmargin=1.2em, itemsep=2pt]
  \item Center: $(-1, 3)$
  \item $a^2 = 16 \Rightarrow a = 4$ (horizontal radius)
  \item $b^2 = 9 \Rightarrow b = 3$ (vertical radius)
\end{itemize}

\columnbreak

\textbf{Step 4 --- Circle as Special Ellipse} (15 min)

When $a = b = r$: $(x-h)^2 + (y-k)^2 = r^2$.

Confirm: $\dfrac{(x-2)^2}{25} + \dfrac{(y+1)^2}{25} = 1$ is a circle,
center $(2,-1)$, $r = 5$.

\textbf{Writing equations:} Center $(3,-2)$, $a = 5$, $b = 2$.
Have students write the ellipse equation, then write a circle equation from
center and radius.

\end{multicols}

\smallskip
\textbf{Pacing:} Reserve 20 min for Tier A activity problems.
\end{skillbox}

\vspace{0.08in}

%----------------------------------------------------------------------
% Lesson — Period 3: Hyperbolas
%----------------------------------------------------------------------
\begin{skillbox}[Lesson (cont.) --- Period 3: Hyperbolas]{sky}
\begin{multicols}{2}

\textbf{Step 5 --- Two Orientations} (12 min)

\begin{itemize}[leftmargin=1.2em, itemsep=2pt]
  \item Opens left/right:
    $\dfrac{(x-h)^2}{a^2} - \dfrac{(y-k)^2}{b^2} = 1$
  \item Opens up/down:
    $\dfrac{(y-k)^2}{b^2} - \dfrac{(x-h)^2}{a^2} = 1$
\end{itemize}

Key diagnostic: the \emph{positive} squared term tells you the direction
of opening.

\columnbreak

\textbf{Step 6 --- Asymptotes} (15 min)

From $\dfrac{(x-2)^2}{9} - \dfrac{(y+1)^2}{4} = 1$:
\begin{itemize}[leftmargin=1.2em, itemsep=2pt]
  \item Center $(2,-1)$; $a = 3$, $b = 2$
  \item Asymptotes: $y + 1 = \pm\dfrac{2}{3}(x - 2)$
\end{itemize}

Contrast with ellipse: same form but minus sign; opens rather than closes;
asymptotes rather than bounded shape.

Students: $\dfrac{(y+3)^2}{16} - \dfrac{(x-1)^2}{9} = 1$ --- identify center,
$a$, $b$, orientation, and asymptotes.

\end{multicols}

\smallskip
\textbf{Pacing:} Exit ticket in final 8 minutes; collect for formative review.
\end{skillbox}

\vspace{0.08in}

%----------------------------------------------------------------------
% Active Monitoring
%----------------------------------------------------------------------
\begin{skillbox}[Active Monitoring]{redbox}
\begin{itemize}[leftmargin=1.4em, itemsep=3pt]
  \item \textbf{Sign of $h$ and $k$}: Students read $(x+1)^2$ and write center
    $(1, \cdot)$ instead of $(-1, \cdot)$. Prompt: \textit{``Rewrite as
    $(x - ?)^2$ so the sign matches the formula.''}
  \item \textbf{Horizontal vs.\ vertical parabola}: Students confuse which
    variable is squared. Prompt: \textit{``Which variable is squared? That
    axis is the axis of symmetry.''}
  \item \textbf{Radii vs.\ $a^2$}: Students report $a = 16$ instead of $a = 4$.
    Prompt: \textit{``The denominator is $a^2$. Take the square root to find $a$.''}
  \item \textbf{Hyperbola direction}: Students pick the wrong orientation.
    Prompt: \textit{``Which term is positive --- the $x$-term or the $y$-term?
    That variable's axis is the direction of opening.''}
  \item \textbf{Asymptote slope}: Students confuse $b/a$ and $a/b$. Prompt:
    \textit{``Write $y - k = \pm(b/a)(x-h)$ and substitute your $a$ and $b$.''}
\end{itemize}
\end{skillbox}

\vspace{0.08in}

%----------------------------------------------------------------------
% Group Work & Differentiation
%----------------------------------------------------------------------
\begin{skillbox}[Group Work \& Differentiation]{redbox}
Students work in groups of 3--4 on the tiered activity (assign by teacher or
allow self-selection with guidance).

\begin{multicols}{3}

\textbf{Tier R --- Remediate}

Given equations in standard form, identify:
\begin{itemize}[leftmargin=1em, itemsep=2pt]
  \item conic type
  \item center or vertex
  \item key measurements ($a$, $b$, or $r$)
\end{itemize}
No conversion or equation-writing required.

\columnbreak

\textbf{Tier A --- Approaching}

Given verbal descriptions (center or vertex and key measurements), write
the standard-form equation. Includes one parabola, one ellipse, and one
hyperbola (with asymptotes).

\columnbreak

\textbf{Tier E --- Extension}

Given general-form equations (e.g.\ $4x^2 - 9y^2 = 36$;
$x^2 + y^2 - 4x + 6y = 3$), rewrite in standard form and identify all
features including asymptotes where applicable.

\end{multicols}
\end{skillbox}

\vspace{0.08in}

%----------------------------------------------------------------------
% Individual Work & Assessment
%----------------------------------------------------------------------
\begin{skillbox}[Individual Work \& Assessment]{redbox}
Exit ticket administered independently (final 8 minutes of Period 3).

\begin{enumerate}[leftmargin=1.4em, itemsep=4pt]
  \item Write the standard form equation of an ellipse centered at $(3, -1)$
    with horizontal radius 5 and vertical radius 2.
  \item Identify the type of conic, its center, and its asymptotes (if any)
    for $\dfrac{(y+2)^2}{16} - \dfrac{(x-1)^2}{9} = 1$.
\end{enumerate}

\textbf{Data use:} Students who cannot identify conic type need reinforcement
reading sign patterns. Students who mix up center signs need targeted reteaching
of the $(x - h)$ convention.
\end{skillbox}

\vspace{0.08in}

%----------------------------------------------------------------------
% Reinforcement & Extension
%----------------------------------------------------------------------
\begin{skillbox}[Reinforcement \& Extension]{goldbox}
\textbf{Homework (Lesson 8.5):} 6 problems covering identifying features from
standard-form equations, writing equations from descriptions, both parabola
orientations, and one AP-style multiple-choice item.

\textbf{Extension:} Given $x^2 + 4y^2 - 6x + 8y - 3 = 0$, complete the square
in both variables to rewrite in standard ellipse form and identify all features.
This bridges to the Tier E activity and previews algebraic techniques used in
calculus.

\textbf{Preview of next lesson:} Ask students: \textit{``We've described conics
with equations in $x$ and $y$. Could you write parametric equations $x(t)$ and
$y(t)$ that trace the same shapes? How might that be useful in modeling
motion?''}
\end{skillbox}

\end{document}
